\section*{الفصل \ref{ch:b1:heat-budget} --- ميزانية حرارة الأرض}

\begin{solution}{exo:b1:heat-budget:1}
$q = k\,\mathrm dT/\mathrm dz = 2.5\times\qty{0.022}{K.m^{-1}} =
\qty{0.055}{W.m^{-2}} = \qty{55}{mW.m^{-2}}$. عبر \qty{1}{km^2} $=
\qty{e6}{m^2}$: \qty{5.5e4}{W}، أي \qty{55}{kW}.
\end{solution}

\begin{solution}{exo:b1:heat-budget:2}
القارات $\num{2.073e14}\times0.0709 = \qty{1.470e13}{W} = \qty{14.70}{TW}$؛
المحيطات $\num{3.028e14}\times0.1054 = \qty{31.92}{TW}$؛ المجموع \qty{46.61}{TW}، و
متوسط \omterm{def:b1:heat-budget:heat-flow}{تدفّق الحرارة} $\num{4.661e13}/\num{5.101e14} = \qty{91.4}{mW.m^{-2}}$ (المنشور 46.7 و\num{91.6} يختلفان فقط بتقريب المتوسطين). ونسبة ما تتلقاه الشمس هي
$46.6/173000 = \num{2.7e-4}$، أي \qty{0.027}{\percent}.
\end{solution}

\begin{solution}{exo:b1:heat-budget:3}
اليورانيوم: $0.99274\times\num{9.513e-5} + 0.00720\times\num{5.685e-4} =
\num{9.444e-5} + \num{4.09e-6} = \qty{9.85e-5}{W.kg^{-1}}$؛ يوفّر النظير النادر
\ce{^{235}U} منه \qty{4}{\percent}. والبوتاسيوم: كسر الكتلة ل\ce{^{40}K} هو $\num{1.17e-4}\times39.96/39.10 = \num{1.196e-4}$، إذن
$\num{1.196e-4}\times\num{2.847e-5} = \qty{3.40e-9}{W.kg^{-1}}$.
\end{solution}

\begin{solution}{exo:b1:heat-budget:4}
$\qty{1}{\micro cal.cm^{-2}.s^{-1}} = \num{4.184e-6}\,\unit{J.s^{-1}}/\qty{e-4}{m^2}
= \qty{41.84}{mW.m^{-2}}$: شرق الولايات المتحدة \qty{33.5}{mW.m^{-2}}، ومنطقة الحوض و
السلاسل \qty{58.6}{mW.m^{-2}}، وسييرا نيفادا \qty{16.7}{mW.m^{-2}}. وب$A_0D$ بوحدة
\unit{mW.m^{-2}} عندما تكون $A_0$ ب\unit{\micro W.m^{-3}} و$D$ ب\unit{km}:
$33.5 + 3.0\times7.5 = \qty{56.0}{mW.m^{-2}}$، $58.6 + 3.0\times9.4 =
\qty{86.8}{mW.m^{-2}}$، $16.7 + 3.0\times10 = \qty{46.7}{mW.m^{-2}}$. ويعطي الغرانيت نفسه تدفّقات حرارة تقريبًا ضعف في منطقة الحوض والسلاسل منها في سييرا نيفادا: فالفرق في \omterm{def:b1:layered-earth:layers}{الوشاح} تحته.
\end{solution}

\begin{solution}{exo:b1:heat-budget:5}
$T(10\,\text{كم}) = 10 + \frac{0.060\times10^4}{2.5} - \frac{\num{2e-6}\times10^8}{2\times2.5}
= 10 + 240 - 40 = \qty{210}{\degreeCelsius}$. تحت ذلك، $q_m = 60 - 2.0\times10 =
\qty{40}{mW.m^{-2}}$، بتدرّج قدره \qty{16}{K.km^{-1}}، و
$T(35\,\text{كم}) = 210 + 16\times25 = \qty{610}{\degreeCelsius}$.
\end{solution}

\begin{solution}{exo:b1:heat-budget:6}
الصخر الطيني: التدرّج $12.0/0.300 = \qty{40}{K.km^{-1}}$، $q = 1.5\times0.040 =
\qty{60}{mW.m^{-2}}$. الغرانيت: $10.0/0.500 = \qty{20}{K.km^{-1}}$،
$q = 3.0\times0.020 = \qty{60}{mW.m^{-2}}$. $R(300) = 300/1.5 =
\qty{200}{m^2.K.W^{-1}}$، $R(800) = 200 + 500/3.0 = \qty{366.7}{m^2.K.W^{-1}}$.
النقاط الثلاث $(0, 12.0)$، $(200, 24.0)$، $(366.7, 34.0)$ على استقامة واحدة، بميل
$22.0/366.7 = \qty{0.0600}{W.m^{-2}}$: \omterm{def:b1:heat-budget:heat-flow}{تدفّق الحرارة} هو نفسه،
\qty{60}{mW.m^{-2}}، في الصخرتين؛ يتغيّر التدرّج وحده، بنسبة عكسية للموصلية.
\end{solution}

\begin{solution}{exo:b1:heat-budget:7}
$\tau = L^2/\kappa$: \qty{1}{m}، \qty{e6}{s} أو \num{12}~يومًا؛ \qty{10}{m}،
\qty{e8}{s} أو \num{3.2}~سنوات؛ \qty{100}{km}، \qty{e16}{s} أو
\qty{317}{Myr}؛ \qty{2900}{km}، \qty{8.4e18}{s} أو \qty{267}{Gyr}. وعلى
\qty{10}{m} لا يُحسّ التغيّر اليومي البتّة، والتغيّر السنوي، الأقصر من السنوات الثلاث اللازمة لعبور \qty{10}{m}، لا يصل إلا ضعيفًا جدًا: لا تصل إلا التغيّرات التي تدوم سنوات أو أكثر. والزمن الأخير هو $267/4.54 = 59$ ضعف عمر الأرض.
\end{solution}

\begin{solution}{exo:b1:heat-budget:8}
(أ) $1646 - 0.54\times50 = \qty{1619}{K}$، أي \qty{1346}{\degreeCelsius}،
وهو المنشور \qty{1350}{\degreeCelsius} مقرَّبًا. (ب) متوسط التدرّج
$(0.54 + 0.36)/2 = \qty{0.45}{K.km^{-1}}$، إذن $1646 + 0.45\times360 =
\qty{1808}{K}$، أي \qty{9}{K} (\qty{0.5}{\percent}) فوق \qty{1799}{K}: التدرّج لا ينخفض خطيًا، لكن التقدير قريب.
\end{solution}

\begin{solution}{exo:b1:heat-budget:9}
(أ) $G^2 = T_0^2/(\pi\kappa t)$ يعطي $t = T_0^2/(\pi\kappa G^2)$.
(ب) $G = 0.02\times\tfrac59\,\unit{K}/\qty{0.3048}{m} = \qty{0.03645}{K.m^{-1}} =
\qty{36.45}{K.km^{-1}}$؛ $t = 3900^2/(\pi\times\num{1.2e-6}\times0.03645^2) =
\num{1.521e7}/\num{5.009e-9} = \qty{3.04e15}{s}$، أي \qty{96}{Myr}، وهو رقم كلفن.
(ج) $t = 1300^2/(\pi\times\num{e-6}\times0.036^2) =
\qty{4.15e14}{s} = \qty{13}{Myr}$: حرارة ابتدائية أكثر واقعية تجعل الأرض الموصِّلة أحدث أيضًا.
\end{solution}

\begin{solution}{exo:b1:heat-budget:10}
(أ) بعد النحت يصبح السطح الجديد هو العمق القديم $e$، حيث كان إنتاج
الحرارة $A_0' = A_0e^{-e/D}$، \omterm{def:b1:heat-budget:heat-flow}{وتدفّق الحرارة} $q(e) = q_m +
A_0De^{-e/D} = q_m + A_0'D$؛ ومقيسًا من السطح الجديد، ما زال الإنتاج يتناقص ك$A_0'e^{-z'/D}$. ويتحرّك الجسم نحو الأسفل على الخط نفسه.
(ب) الطبقة صارت الآن سمكها $h - e$ بال$A$ نفسه، إذن $q_0' = q_m + A(h - e)$:
على خط ميله $h - e$ يمرّ ب$q_m$. (ج) النموذج الأسّي وحده: ففي نموذج الطبقة المنتظمة، ستقع أجسام نُحتت بأعماق مختلفة على خطوط بميول مختلفة، وتتبعثر نقاط منطقة واحدة. وأن يُلاحظ خط واحد يرجّح إنتاجًا يتناقص تقريبًا أسّيًا مع العمق.
\end{solution}

\begin{solution}{exo:b1:heat-budget:11}
(أ) $\tau_P = RL/(3\kappa)$: من أجل \qty{50}{km}، $\num{6.4e6}\times\num{5e4}/\num{3e-6}
= \qty{1.07e17}{s} = \qty{3.4}{Gyr}$؛ ومن أجل \qty{100}{km}، \qty{6.8}{Gyr}.
(ب) $q = kT_i/L = 3\times1300/\num{5e4} = \qty{78}{mW.m^{-2}}$ في البداية،
و$78\,e^{-2/3.38} = \qty{43}{mW.m^{-2}}$ بعد \qty{2}{Gyr}. (ج)
$3\times1300/\sqrt{\pi\times\num{e-6}\times\num{6.31e16}} = 3900/\num{4.45e5} =
\qty{8.8}{mW.m^{-2}}$: فالأرض الممزوجة جيدًا ما زالت تفقد حرارة أكثر بخمسة أضعاف بعد مليارَي سنة. فكامل داخلها يغذّي السطح، بينما لا يفعل ذلك في أرض موصِّلة إلا \omterm{def:b1:layered-earth:layers}{قشرة} خارجية.
\end{solution}

\begin{solution}{exo:b1:heat-budget:12}
(أ) $\sqrt{\pi\kappa\times\qty{1}{Myr}} = \sqrt{\pi\times\num{e-6}\times\num{3.156e13}}
= \qty{9957}{m}$، إذن $C = 3\times1336/9957 = \qty{0.403}{W.m^{-2}}$ من أجل
$t$ بال Myr، \qty{403}{mW.m^{-2}.Myr^{1/2}}: \qty{82}{\percent} من ال\num{490} المطابَقة. ويحتاج النموذج إلى $k\,\Delta T$ أكبر (نحو \qty{4900}{W.m^{-1}}): موصلية أعلى \omterm{def:b1:layered-earth:layers}{للوشاح}، أو داخل أسخن، كما على المنحنى الأديباتي.
(ب) $490/\sqrt1 = 490$، $490/\sqrt{10} = 155$ و$490/\sqrt{100} =
\qty{49}{mW.m^{-2}}$. (ج) $z = 2.326\sqrt{\num{e-6}\times\num{3.156e15}} =
2.326\times\qty{56180}{m} = \qty{131}{km}$.
\end{solution}

\begin{solution}{pb:b1:heat-budget:1}
\textbf{1.} التدرّج $= (T_{500} - T_{100})/\qty{0.4}{km}$: A \qty{15.25}{K.km^{-1}}،
B \num{19.50}، C \num{19.50}، D \num{25.50}، E \qty{28.25}{K.km^{-1}}.

\textbf{2.} يُخلّ الحفر بحرارة الصخر (يبرّد سائل الحفر جدران الثقب أو يسخّنها)، ويستغرق زوال الاضطراب أسابيع إلى شهور. أما أعلى عشرات الأمتار فتتبع الدورة الفصلية والتغيّرات الحديثة لحرارة السطح.

\textbf{3.} $q_0 = k\times$ التدرّج: A $2.9\times15.25 = \qty{44.225}{mW.m^{-2}}$،
B \qty{50.70}{mW.m^{-2}}، C \qty{60.45}{mW.m^{-2}}، D \qty{68.85}{mW.m^{-2}}،
E \qty{84.75}{mW.m^{-2}}.

\textbf{4.} A إلى D أدنى من المتوسط القاري، وE أعلى منه ب\qty{13.9}{mW.m^{-2}}، أي \qty{20}{\percent}.

\textbf{5.} \omterm{def:b1:heat-budget:heat-flow}{تدفّق الحرارة} هو الموصلية مضروبة في التدرّج: C يوصّل أفضل (3.1 مقابل 2.6)، فالتدرّج نفسه ينقل \qty{19}{\percent} حرارة أكثر.

\textbf{6.} تتبع متساويات الحرارة التضاريس جزئيًا فقط: تحت القمة ترتفع أقل من السطح، فتتباعد تحتها، فيقلّ التدرّج المقاس (ويزيد تحت الوادي). ويلزم تصحيح للتضاريس قبل مقارنة الأجسام.

\textbf{7.} $\rho C_U H_U$: A $2650\times\num{2.0e-6}\times\num{9.85e-5} =
\qty{0.522}{\micro W.m^{-3}}$، B 0.914، C 1.305، D 1.827، E 2.610.

\textbf{8.} الثوريوم: A 0.558، B 0.976، C 1.255، D 1.742، E 2.439.

\textbf{9.} البوتاسيوم: A $2650\times0.030\times\num{3.40e-9} = 0.270$، B 0.315،
C 0.360، D 0.378، E 0.405.

\textbf{10.} $A_0$ = 1.35، 2.20، 2.92، 3.95 و\qty{5.46}{\micro W.m^{-3}}.
يسهم البوتاسيوم بأقل قدر، من \qty{20}{\percent} من المجموع في A إلى \qty{7}{\percent} في E؛ ويتقاسم اليورانيوم والثوريوم الباقي بالتساوي تقريبًا.

\textbf{11.} نعم: النقاط الخمس ترتفع بانتظام، ضمن نحو \qty{1.5}{mW.m^{-2}} من خط مستقيم.

\textbf{12.} $\bar A_0 = 15.88/5 = 3.176$، $\bar q_0 = 308.975/5 = 61.795$؛
$\sum(A_0 - \bar A_0)^2 = 10.1681$؛ $\sum(A_0 - \bar A_0)(q_0 - \bar q_0) =
101.146$.

\textbf{13.} $D = 101.146/10.1681 = 9.95$ (بوحدة \unit{mW.m^{-2}} لكل
\unit{\micro W.m^{-3}}، أي بالكيلومتر): $D = \qty{9.95}{km}$، و$q_m = 61.795 -
9.95\times3.176 = \qty{30.20}{mW.m^{-2}}$.

\textbf{14.} الميل، نحو \qty{10}{km}، هو ميل المقاطعات الثلاث؛ \omterm{def:b1:heat-budget:heat-flow}{وتدفّق الحرارة} المخفّض، \qty{30}{mW.m^{-2}}، قريب من \qty{33}{mW.m^{-2}} لشرق الولايات المتحدة: مقاطعة هادئة ذات \omterm{def:b1:layered-earth:layers}{وشاح} بارد تحتها، بخلاف منطقة الحوض والسلاسل.

\textbf{15.} $q_m/k = \qty{12}{K.km^{-1}}$ و$A_0D^2/k = \num{5.46e-6}\times\num{e8}/2.5
= \qty{218.4}{K}$: $T(z) = 10 + 12z + 218.4\,(1 - e^{-z/10})$، $z$ بالكيلومتر، $T$ ب\unit{\degreeCelsius}.

\textbf{16.} $q(z) = 30 + 54.6\,e^{-z/10}$: \qty{50.1}{mW.m^{-2}} عند \qty{10}{km}،
\num{37.4} عند \qty{20}{km} و\qty{31.7}{mW.m^{-2}} عند \qty{35}{km}.

\textbf{17.} $30/84.6 = \qty{35}{\percent}$: نحو ثلثي \omterm{def:b1:heat-budget:heat-flow}{تدفّق الحرارة} فوق E ينتج في طبقته المشعّة.

\textbf{18.} $T(10) = 10 + 120 + 218.4\times0.632 = \qty{268}{\degreeCelsius}$.

\textbf{19.} $T(20) = 10 + 240 + 218.4\times0.865 = \qty{439}{\degreeCelsius}$.

\textbf{20.} تحت A، $A_0D^2/k = 1.35\times100/2.5 = \qty{54.0}{K}$:
$T(35) = 10 + 420 + 54.0\times0.970 = \qty{482}{\degreeCelsius}$.

\textbf{21.} ب$A_0D^2/k$: \qty{54}{K} تحت A و\qty{218}{K} تحت E. \omterm{def:b1:heat-budget:heat-flow}{تدفّق الحرارة} العميق هو نفسه؛ \omterm{def:b1:layered-earth:layers}{والقشرة} في E أغنى بمصادر الحرارة أربع مرات، والحرارة التي تنتجها يجب أن توصَّل صعودًا عبر الصخر نفسه، وهذا يحتاج إلى تدرّج أشدّ طوال الطريق.

\textbf{22.} $1300 = 10 + 218.4 + 12z$ يعطي $z = \qty{89.3}{km}$ تحت E؛
$1300 = 10 + 54.0 + 12z$ يعطي \qty{103.0}{km} تحت A.

\textbf{23.} يقع كلاهما ضمن مدى \omterm{def:b1:layered-earth:lithosphere}{الغلاف الصخري} القاري، من مئة كيلومتر أو أقل تحت المناطق الفتية إلى نحو \qty{250}{km} تحت الكراتونات، نحو طرفه الرقيق: فمع تدفّق حرارة مخفّض قدره \qty{30}{mW.m^{-2}} وهذه الموصلية، يبلغ \omterm{def:b1:heat-budget:geotherm}{المنحنى الحراري} الموصل حرارة \omterm{def:b1:layered-earth:layers}{الوشاح} عند نحو \qty{100}{km}؛ وتحت الجسم الأكثر إشعاعًا، أبكر ب\qty{14}{km}.

\textbf{24.} $T(35) = 10 + 420 + 218.4\times(1 - e^{-3.5}) = 430 + 211.8$:
\omterm{def:b1:layered-earth:discontinuity}{موهو} تحت الجسم E على \textbf{\qty[text-series-to-math=true]{642}{\degreeCelsius}}.
\end{solution}
